\section{Conclusions}
\label{sec:conclusion}

We have modelled the human urinary bladder as a fluid-filled viscoelastic
spherical shell and computed its flexural mode frequencies across the
physiological fill range.  The principal findings are:

\begin{enumerate}
  \item The fundamental flexural mode ($n = 2$) lies between \SI{14}{\hertz}
    and \SI{17}{\hertz} for fill volumes of \SIrange{150}{500}{\milli\litre}
    where the thin-shell assumption is valid ($h/R < 0.1$).
    A frequency minimum of approximately \SI{14}{\hertz} occurs near
    \SI{220}{\milli\litre}, arising from the competition between geometric
    softening (increasing radius and thinning wall) and strain-stiffening of
    the detrusor muscle.

  \item Mechanical vibration transmitted through the pelvic skeleton is
    approximately $6{,}500$ times more effective than airborne sound at
    exciting the bladder's $n = 2$ mode at \SI{300}{\milli\litre} fill.
    This coupling disparity, driven by the $(ka)^2$ penalty in the Rayleigh
    regime ($ka \approx 0.01$), confirms that whole-body vibration is the
    dominant excitation pathway---consistent with the epidemiological
    evidence and with the analogous result for the abdominal cavity.

  \item The bladder's predicted resonance range (\SIrange{14}{17}{\hertz})
    sits above the ISO~2631 peak pelvic transmissibility band
    (\SIrange{4}{8}{\hertz}).  At typical occupational WBV frequencies, the
    bladder therefore responds as a forced sub-resonant oscillator driven by
    amplified pelvic motion, rather than being excited at its own natural
    frequency.

  \item The breathing mode ($f_0 \approx \SI{9500}{\hertz}$) is dominated by
    the fluid bulk modulus and is irrelevant to WBV or infrasonic exposure,
    confirming the same qualitative separation between breathing and flexural
    mode families observed in the larger abdominal cavity.
\end{enumerate}

The model provides a first-principles basis for understanding
vibration-induced bladder symptoms.  Under representative occupational WBV
exposure ($a = \SI{0.5}{\metre\per\second\squared}$), the predicted peak
bladder wall displacement is approximately \SI{55}{\micro\metre} at resonance
and \SI{170}{\micro\metre} in the sub-resonant regime---five orders of
magnitude smaller than normal filling displacements, yet delivered at cyclic
strain rates ($\sim$\SI{0.1}{\per\second}) that are 500 times faster than
quasi-static filling and thus relevant to the rate-sensitive PIEZO2
mechanotransduction pathway.  These results may inform both occupational
exposure guidelines and the design of therapeutic WBV protocols for pelvic
floor rehabilitation.  Future work should characterise the dynamic activation
thresholds of bladder mechanoreceptors at infrasonic frequencies, incorporate
the pelvic boundary constraint through finite element analysis with
anatomically realistic geometry, and pursue experimental validation of the
predicted resonance range.
